% docs/texte/90_appendix/figures/fringe_pattern.tex -- fringe pattern of a two-element interferometer (pgfplots)
% Version: 0.1.1 (2026-10-12) -- legend above the plot
% Normalised response R(theta) = cos(2 pi (b/lambda) cos theta) for b = 10 lambda,
% monochromatic and with a fractional bandwidth of 10 % (envelope
% |sinc(dnu tau_g)|, dnu tau_g = (dnu/nu)(b/lambda) cos theta).
\begin{tikzpicture}
  \begin{axis}[
      width=0.95\textwidth, height=6cm,
      xmin=0, xmax=180, ymin=-1.15, ymax=1.15,
      xtick={0,30,60,90,120,150,180},
      xlabel={angle $\theta$ between baseline and source direction [deg]},
      ylabel={normalised response $R/R_{\max}$},
      legend style={font=\footnotesize, at={(0.5,1.12)}, anchor=south, draw=none},
      legend columns=3,
      tick label style={font=\footnotesize}, label style={font=\footnotesize},
      samples=1200, domain=0:180]
    % monochromatic
    \addplot[thin, gray] {cos(deg(2*pi*10*cos(x)))};
    \addlegendentry{monochromatic}
    % 10 % bandwidth: fringe times sinc envelope
    \addplot[thick, blue!70!black]
      {cos(deg(2*pi*10*cos(x))) * sin(deg(pi*0.1*10*cos(x)+1e-9))/(pi*0.1*10*cos(x)+1e-9)};
    \addlegendentry{$\Delta\nu/\nu = 0.1$}
    \addplot[thin, dashed, red!70!black]
      {abs(sin(deg(pi*0.1*10*cos(x)+1e-9))/(pi*0.1*10*cos(x)+1e-9))};
    \addlegendentry{envelope $|\operatorname{sinc}(\Delta\nu\,\tau_g)|$}
    % mark fringe spacing near theta = 90 deg
    \draw[<->, thick] (axis cs:90,1.07) -- (axis cs:95.73,1.07)
      node[midway, above, font=\scriptsize] {$\approx\lambda/b$};
  \end{axis}
\end{tikzpicture}
